\chapter*{Resumen}


El trabajo aborda el diseño desarrollo e implementación de una red social orientada a compartir y distribuir contenido audiovisual en línea.  El objetivo principal es desarrollar un prototipo funcional de una plataforma de streaming de videos, denominada Booster, que permita a los usuarios publicar, compartir, descubrir e interactuar con contenido digital de manera sencilla e intuitiva centrada en la comunidad.

Booster surge así como respuesta a problemas existentes identificados por múltiples usuarios en las plataformas de videos actuales. En particular la limitada interacción entre creadores y audiencias; los algoritmos de recomendación opacos; la escasa participación de los usuarios en el proceso de descubrimiento de comunidades; la gran cantidad de contenido publicitario; la creciente creación de contenido con base en IA sin la debida advertencia y la ausencia de mecanismos que incentiven la contribución dentro de una comunidad, solo por mencionar algunos. Para resolver estos problemas, Booster introduce elementos de gamificación y mecanismos de interacción social para fomentar la participación, fortalecer el sentido de pertenencia de la comunidad y reducir el contenido comercial.

Adicionalmente, el trabajo realiza un análisis completo de la arquitectura y las tecnologías disponibles.  Se evalúan distintos proveedores de infraestructura en la nube, servicios dedicados al procesamiento y distribución de video, sistemas de almacenamiento, bases de datos y herramientas de desarrollo, con el fin de alinear las decisiones de diseño de forma eficiente con los objetivos técnicos y empresariales del proyecto.

El trabajo recopila  el proceso completo de ingeniería de software desde la definición de requisitos y el diseño de la arquitectura hasta la implementación de funcionalidades y la puesta en marcha del servicio en un entorno de producción.  De ese modo,  la solución busca ofrecer una alternativa a las plataformas de streaming de video tradicionales que pretende optimizar el descubrimiento de contenido,  aumentar las interacciones sociales y la participación en comunidades. 

Finalmente,  el trabajo demuestra a través del desarrollo de un Mínimo Producto Viable (MVP) \href{http://boostervideos.net}{\underline{www.boostervideos.net}} la viabilidad técnica de una plataforma de este tipo que introduce nuevas funcionalidades más allá de la simple distribución de contenido en busca de una experiencia de usuario más dinámica,  orientada a la participación y a la creación de comunidades digitales.


% El trabajo aborda el diseño, desarrollo e implementación de una red social orientada a compartir y consumir contenido audiovisual en línea. El objetivo principal es desarrollar un prototipo funcional de una plataforma de streaming de videos, llamada \textit{Booster}, que permita a los usuarios publicar, compartir, descubrir e interactuar con contenido digital de manera sencilla e intuitiva, centrada en la comunidad.

% La propuesta surge como respuesta a problemas existentes en las plataformas de videos actuales. Entre estos, se encuentra la limitada interacción entre creadores y audiencias; los algoritmos de recomendación opacos; la escasa participación de los usuarios en el proceso de descubrimiento de comunidades; la gran cantidad de contenido publicitario y, la ausencia de mecanismos que incentiven la contribución dentro de una comunidad. Frente a estos problemas, \textit{Booster} emplea elementos mecanismos de gamificación y elementos de interacción social para fomentar la participación, fortalecer el sentido de pertenencia a la comunidad y reducir el contenido comercial. 

% Adicionalmente, se realiza un análisis completo de la arquitectura y las tecnologías disponibles. Se evalúan distintos proveedores de infraestructura en la nube, servicios dedicados al procesamiento y distribución de video, sistemas de almacenamiento, bases de datos y herramientas de desarrollo, con el fin de alinear las decisiones de diseño de la forma eficiente con los objetivos técnicos y empresariales del proyecto.

% El trabajo recopila el proceso completo de ingeniería de software desde la definición de requisitos y el diseño de la arquitectura hasta la implementación de funcionalidades y la puesta en marcha del servicio en un entorno de producción. De este modo, la solución busca ofrecer una alternativa a las plataformas de streaming de video tradicionales que pretende optimizar el descubrimiento de contenido, aumentar las interacciones sociales y la participación en comunidades. 

% En definitiva, el trabajo demuestra la viabilidad técnica de una plataforma de este tipo que introduce nuevas funcionalidades aparte de la simple distribución de contenido, con el objetivo de obtener una experiencia de usuario más dinámica, orientada a la participación y a la creación de comunidades digitales. 






% El trabajo aborda el diseño y la implementación de una red social de vídeos en formato web llamada Booster. Se detallará el funcionamiento de la infraestructura de la plataforma así como las tecnologías utilizadas. Además, se hará una evaluación de las distintas opciones para procesar los vídeos: Bunny.net, Mux y AWS. También se documentará el código, las decisiones de diseño y se realizará una guía de despliegue. 

% El objetivo del proyecto es implementar un prototipo (Producto Mínimo Viable) y detallar cómo funcionan todas las partes de la plataforma de vídeos: 

% \begin{itemize}
%     \item autenticación
%     \item arquitectura del procesamiento y transcodificación de vídeos
%     \item  despliegue,
%     \item diseños de las interfaces de usuario
%     \item optimizaciones
%     \item riesgos de seguridad
%     \item creación de comentarios
%     \item notificaciones
%     \item diseño del algoritmo de recomendación
%     \item creación de comunidades
%     \item carga de vídeos y edición de sus metadatos
%     \item recuento de visitas
%     \item calificaciones a los vídeos
%     \item filtrado de los vídeos por categorías
%     \item recuento de los seguidores de un usuario
%     \item explicación del sistema de Boost y cómo funciona la moneda virtual (XP)
%     \item diseño de la base de datos
%     \item conexión entre front-end y back-end y explicación de las tecnologías usadas:

% \begin{itemize}
% \item Next.Js
% \item tRPC
% \item NeonDB
% \item DrizzleORM
% \item SupaBase
% \item TailwindCSS
% \item ReactJS
% \end{itemize}
% \end{itemize}
\begin{otherlanguage}{english}
  \chapter*{Abstract}

This work addresses the design, development, and implementation of a social network for sharing and distributing online audiovisual content. The main objective is to develop a functional prototype of a video streaming platform, called Booster, that allows users to easily and intuitively publish, share, discover, and interact with digital content in a community-centric way.

Booster thus emerges as a response to existing problems identified by many users on current video platforms. These include limited interaction between creators and audiences; opaque recommendation algorithms; low user participation in the community discovery process; the large amount of advertising content; the increasing creation of AI-based content without proper warning; and the lack of mechanisms to incentivize contributions within a community, to name just a few. To solve these problems, Booster introduces gamification elements and social interaction mechanisms to encourage participation, strengthen community sense of belonging, and reduce commercial content.

Additionally, this work includes a comprehensive analysis of the architecture and available technologies. Various cloud infrastructure providers, dedicated video processing and distribution services, storage systems, databases, and development tools are evaluated to efficiently align design decisions with the project's technical and business objectives.

This work compiles the complete software engineering process, from requirements definition and architecture design to functional implementation and service deployment in a production environment. The solution aims to offer an alternative to traditional video streaming platforms, optimizing content discovery, increasing social interactions, and fostering community engagement.

Finally, the work demonstrates, through the development of a Minimum Viable Product (MVP) \href{http://boostervideos.net}{\underline{www.boostervideos.net}}, the technical feasibility of such a platform, which introduces new functionalities beyond simple content distribution, striving for a more dynamic user experience focused on engagement and the creation of digital communities.
  
\end{otherlanguage}


%%%%%%%%%%%%%%%%%%%%%%%%%%%%%%%%%%%%%%%%%%%%%%%%%%%%%%%%%%%
%% Final del resumen.
%%%%%%%%%%%%%%%%%%%%%%%%%%%%%%%%%%%%%%%%%%%%%%%%%%%%%%%%%%%